% GENERATED FILE. Edit canonical lessons, then run npm run generate:books.
% canonical-source-hash: fnv1a64:c841117bed603d91
% canonical-lessons: ML-C78-nre, ML-C78-ude, ML-R78-genitive-recall

\chapter{Whose It Is}
\label{ch:ml-genitive}
\hlchaptermodality{78}

\begin{chapteropening}
\textbf{By the end of this chapter:} \emph{I can say that any noun owns something, not only that I do.}

From cold: both genitive endings, what decides between them, and the owner-first order the phrase always takes.
\end{chapteropening}

\section[adhyāpakanṟe]{\ml{അധ്യാപകന്റെ} (adhyāpakanṟe) — \textquotedblleft{}the teacher's\textquotedblright{}}

\label{lesson:ML-C78-nre}



\begin{quote}
Say \emph{my}. Then say the word for \emph{teacher}.

You have owned the first of those since the second chapter, and you have never been able to do to any other word what it does.
\end{quote}

\begin{grammarlens}[title={Grammar Lens: the ending was there all along}]
\textbf{\ml{എന്റെ}} was handed to you whole, as a single word meaning \emph{my}. Look at its end:

\noindent
\begin{tabularx}{\linewidth}{@{}>{\raggedright\arraybackslash}X>{\raggedright\arraybackslash}X@{}}
\toprule
\textbf{} & \textbf{} \\
\midrule
\ml{എ}\textbf{\ml{ന്റെ}} & my \\
\ml{അധ്യാപക}\textbf{\ml{ന്റെ}} & the teacher's \\
\bottomrule
\end{tabularx}

\textbf{-\ml{ന്റെ} is a suffix, and it makes whatever it lands on the owner.} That is the one thing this lesson adds, and it turns a word you memorised into a machine you can run.

\begin{quote}
\textbf{\ml{അധ്യാപകന്റെ} \ml{പുസ്തകം}}
\end{quote}

\begin{quote}
\emph{adhyāpakanṟe pustakaṁ} — \textquotedblleft{}\textbf{the teacher's book}\textquotedblright{}
\end{quote}

The owner comes first, wearing the ending; the thing owned follows bare. English does the same with \emph{'s} and the opposite with \emph{of}. Malayalam only has the first order.
\end{grammarlens}

\begin{cousinweb}
The script lesson for \textbf{\ml{എന്റെ}} already showed you \textbf{\ml{ന്റെ}} on its own, as a conjunct to read. It was a piece of writing then. It is a piece of grammar now, and it is the same piece.

That is worth expecting from this language: \textbf{something introduced as a shape often turns out later to be a part.}
\end{cousinweb}

\subsection*{Your turn}
\begin{itemize}
\raggedright
  \item \emph{Say it:} \emph{adhyāpakanṟe} — the teacher's
  \item \emph{Say it:} \emph{adhyāpakanṟe pustakaṁ}
  \item \emph{Look:} at \textbf{\ml{എന്റെ}} and \textbf{\ml{അധ്യാപകന്റെ}}, and find where they agree
  \item \emph{Say it:} which of the two words in the phrase is the owner
\end{itemize}

\begin{tcolorbox}[breakable,colback=teal!4,colframe=teal!35!black,title={Before you move on}]
What does \textbf{-\ml{ന്റെ}} do to a noun? (\textbf{Makes it the owner}.) How do you say \emph{the teacher's book}? (\emph{\textbf{\ml{അധ്യാപകന്റെ} \ml{പുസ്തകം}}}.) And which comes first, the owner or the thing? (\textbf{The owner}.)
\end{tcolorbox}
\section[ammayuṭe]{\ml{അമ്മയുടെ} (ammayuṭe) — \textquotedblleft{}mother's\textquotedblright{}}

\label{lesson:ML-C78-ude}



\begin{quote}
Say \emph{the teacher's book}. Then take the ending off and say the teacher again on their own.
\end{quote}

\begin{grammarlens}[title={Grammar Lens: the same job, a different ending}]
\begin{quote}
\textbf{\ml{അമ്മയുടെ} \ml{പുസ്തകം}}
\end{quote}

\begin{quote}
\emph{ammayuṭe pustakaṁ} — \textquotedblleft{}\textbf{mother's book}\textquotedblright{}
\end{quote}

Not \textbf{-\ml{ന്റെ}}. \textbf{-\ml{ഉടെ}}, and a \textbf{\ml{യ}} appears in front of it.

\textbf{\ml{അധ്യാപകൻ} takes -\ml{ന്റെ}. \ml{അമ്മ} takes -\ml{യുടെ}.} Two nouns, two landings.

\noindent
\begin{tabularx}{\linewidth}{@{}>{\raggedright\arraybackslash}X>{\raggedright\arraybackslash}X@{}}
\toprule
\textbf{this noun} & \textbf{ending} \\
\midrule
\textbf{\ml{അധ്യാപകൻ}} & \textbf{-\ml{ന്റെ}} \\
\textbf{\ml{അമ്മ}} & \textbf{-\ml{യുടെ}} \\
\bottomrule
\end{tabularx}

\textbf{Which of the two a noun takes is learned with the noun}, not worked out from its last sound — and you have had the proof since your first chapter. \textbf{\ml{നിന്റെ}} is the owner-form of \textbf{\ml{നീ}}. \textbf{\ml{നീ}} ends in a vowel, exactly as \textbf{\ml{അമ്മ}} does, and it does not take \textbf{-\ml{യുടെ}} at all.

\textbf{What is reliable is the shape of the phrase.} Owner first with its ending, then the thing owned — that never varies, and it is what lets any noun in this book own something.
\end{grammarlens}

\subsection*{Your turn}
Say these aloud:

\begin{itemize}
\raggedright
  \item \emph{ammayuṭe} — mother's
  \item \emph{ammayuṭe pustakaṁ}
  \item both owners in a row — \emph{adhyāpakanṟe}, \emph{ammayuṭe}
  \item what decided which ending each one took
\end{itemize}

\begin{tcolorbox}[breakable,colback=teal!4,colframe=teal!35!black,title={Before you move on}]
Which ending does \textbf{\ml{അമ്മ}} take? (\textbf{-\ml{യുടെ}}.) How do you say \emph{mother's book}? (\emph{\textbf{\ml{അമ്മയുടെ} \ml{പുസ്തകം}}}.) And who picks the ending — you, or the noun? (\textbf{The noun}.)
\end{tcolorbox}
\section[ārute enna ōrmma]{(whose it is) — from cold}

\label{lesson:ML-R78-genitive-recall}



\begin{quote}
Close the lessons before this one. Say \emph{my}, and then say what its last piece does when you put it on somebody else.
\end{quote}

\begin{grammarlens}[title={Grammar Lens: two endings, one job, and the noun decides}]
\noindent
\begin{tabularx}{\linewidth}{@{}>{\raggedright\arraybackslash}X>{\raggedright\arraybackslash}X@{}}
\toprule
\textbf{owner} & \textbf{ending} \\
\midrule
\textbf{\ml{അധ്യാപകന്റെ}} & \textbf{-\ml{ന്റെ}} \\
\textbf{\ml{അമ്മയുടെ}} & \textbf{-\ml{യുടെ}} \\
\bottomrule
\end{tabularx}

One job — \emph{whose it is} — and \textbf{which of the two endings a noun takes comes with the noun.} \textbf{\ml{നിന്റെ}} already showed you that: \textbf{\ml{നീ}} ends in a vowel, as \textbf{\ml{അമ്മ}} does, and takes \textbf{-\ml{ന്റെ}} all the same.

Then the phrase is built the same way every time: \textbf{owner first with the ending, thing second and bare.} \emph{adhyāpakanṟe pustakaṁ}, \emph{ammayuṭe pustakaṁ}.
\end{grammarlens}

\begin{grammarlens}[title={What you've built}]
A third possessor, and every one after it.

\textbf{\ml{എന്റെ}} and \textbf{\ml{നിന്റെ}} were handed over whole — \emph{my} and \emph{your}, two words to memorise. Naming the ending inside them turns a closed pair into an open rule: \textbf{any noun this book has taught you can now own something.}

That is the difference between a word and a piece of grammar, and it is worth looking for: when two words you were given separately end the same way, the ending is usually doing the work.
\end{grammarlens}

\subsection*{Your turn}
Say these aloud:

\begin{itemize}
\raggedright
  \item \emph{adhyāpakanṟe pustakaṁ}
  \item \emph{ammayuṭe pustakaṁ}
  \item which ending each took, and what decided it
  \item \emph{my book}, using the word you have had since the second chapter
  \item which comes first in the phrase, the owner or the thing
\end{itemize}

\begin{tcolorbox}[breakable,colback=teal!4,colframe=teal!35!black,title={Before you move on}]
Which ending does \textbf{\ml{അധ്യാപകൻ}} take? (\textbf{-\ml{ന്റെ}}.) And \textbf{\ml{അമ്മ}}? (\textbf{-\ml{യുടെ}}.) Can you tell which one a new noun will take by listening to its end? (\textbf{No} — \textbf{\ml{നീ}} ends like \textbf{\ml{അമ്മ}} and takes \textbf{\ml{നിന്റെ}}.) And what does either one do to the noun it lands on? (\textbf{Makes it the owner}.)
\end{tcolorbox}
